\documentclass[11pt,a4paper]{article}
\usepackage[utf8]{inputenc}
\usepackage{amsmath,amssymb,amsthm}
\usepackage{listings}
\usepackage[margin=2.5cm]{geometry}
\usepackage{xcolor}

\newtheorem{theorem}{Theorem}
\newtheorem{lemma}[theorem]{Lemma}

\lstset{
  language=C++,
  basicstyle=\ttfamily\small,
  keywordstyle=\color{blue}\bfseries,
  commentstyle=\color{green!50!black},
  stringstyle=\color{red!70!black},
  numbers=left,
  numberstyle=\tiny\color{gray},
  breaklines=true,
  frame=single,
  tabsize=4,
  showstringspaces=false
}

\title{IOI 1990 -- Problem 3: Festival Decoration}
\author{Solution}
\date{}

\begin{document}
\maketitle

\section{Problem Statement}

Given $k$ colors with counts $c_1, c_2, \ldots, c_k$ (totaling
$n = c_1 + c_2 + \cdots + c_k$ balls), arrange all $n$ balls in a line such
that no two adjacent balls share the same color. Output one valid arrangement,
or report that none exists.

\section{Solution}

\subsection{Feasibility Condition}

\begin{theorem}
A valid arrangement exists if and only if
\[
  \max_i\, c_i \le \left\lceil \frac{n}{2} \right\rceil.
\]
\end{theorem}
\begin{proof}
\textbf{Necessity.} In any valid arrangement of $n$ elements, a single color
can occupy at most every other position, i.e., at most $\lceil n/2 \rceil$
positions.

\textbf{Sufficiency.} The greedy algorithm below always produces a valid
arrangement when this condition holds, as shown in the correctness argument
that follows.
\end{proof}

\subsection{Greedy Algorithm}

Maintain a max-heap of $(\text{count}, \text{color})$ pairs:

\begin{enumerate}
  \item Extract the color with the highest remaining count.
  \item If it differs from the previously placed color, place it.
  \item Otherwise, extract the second-highest color, place it, and push
    the first color back.
  \item Repeat until all balls are placed.
\end{enumerate}

\begin{lemma}[Correctness]
If $\max_i\, c_i \le \lceil n/2 \rceil$, the greedy algorithm never gets
stuck.
\end{lemma}
\begin{proof}
The algorithm gets stuck only when the heap contains a single color that
matches the last placed color. At that point, the remaining count of that
color must be at least~1 while the total remaining is also that count. But
this means $c_{\max} > \lceil n'/2 \rceil$ for the remaining $n'$ items,
which contradicts the invariant maintained by always placing the most frequent
available color first.
\end{proof}

\section{Complexity Analysis}

\begin{itemize}
  \item \textbf{Time:} $O(n \log k)$. Each of the $n$ placement steps
    involves at most two heap operations, each costing $O(\log k)$.
  \item \textbf{Space:} $O(k)$ for the heap.
\end{itemize}

\section{Implementation}

\begin{lstlisting}
#include <iostream>
#include <queue>
#include <string>
using namespace std;

int main() {
    int k;
    cin >> k;

    priority_queue<pair<int,int>> pq;
    int total = 0;

    for (int i = 0; i < k; i++) {
        int c;
        cin >> c;
        if (c > 0) {
            pq.push({c, i});
            total += c;
        }
    }

    if (pq.empty()) {
        cout << endl;
        return 0;
    }

    if (pq.top().first > (total + 1) / 2) {
        cout << "IMPOSSIBLE" << endl;
        return 0;
    }

    string result;
    int lastColor = -1;

    while (!pq.empty()) {
        auto [cnt1, col1] = pq.top();
        pq.pop();

        if (col1 != lastColor) {
            result += ('A' + col1);
            lastColor = col1;
            if (cnt1 - 1 > 0)
                pq.push({cnt1 - 1, col1});
        } else {
            if (pq.empty()) {
                // Should not happen if feasibility check passed
                cout << "IMPOSSIBLE" << endl;
                return 0;
            }
            auto [cnt2, col2] = pq.top();
            pq.pop();

            result += ('A' + col2);
            lastColor = col2;

            pq.push({cnt1, col1});
            if (cnt2 - 1 > 0)
                pq.push({cnt2 - 1, col2});
        }
    }

    cout << result << endl;
    return 0;
}
\end{lstlisting}

\section{Example}

\textbf{Input:} 3 colors with counts $[3, 2, 1]$, total $n = 6$.

Since $\max(c_i) = 3 = \lceil 6/2 \rceil$, a valid arrangement exists.

\medskip
\noindent Greedy execution:
\begin{enumerate}
  \item Heap: $\{(3,\mathrm{A}),\, (2,\mathrm{B}),\, (1,\mathrm{C})\}$.
    Place \texttt{A}. Remaining: $[2,2,1]$.
  \item Top is $\mathrm{A}$ again (tied with $\mathrm{B}$); last was
    $\mathrm{A}$, so take $\mathrm{B}$. Place \texttt{B}. Remaining:
    $[2,1,1]$.
  \item Place \texttt{A}. Remaining: $[1,1,1]$.
  \item Place \texttt{B}. Remaining: $[1,0,1]$.
  \item Place \texttt{A}. Remaining: $[0,0,1]$.
  \item Place \texttt{C}. Done.
\end{enumerate}

Result: \texttt{ABABAC}.

\end{document}
